\documentclass[11pt,a4paper]{article}
\usepackage[margin=23mm]{geometry}
\usepackage{amsmath,amssymb,amsthm}
\setlength{\emergencystretch}{2em}
\newtheorem{lemma}{Lemma}
\newtheorem{theorem}[lemma]{Theorem}
\newcommand{\gen}[1]{\langle#1\rangle}
\title{Choosing the local witnesses simultaneously}
\author{Human Compression project: L044 continuation}
\date{}
\begin{document}\maketitle
All marks below are involutions. Retain every actual
odd-degree head and global paired-tail hypothesis of
\emph{Local pairs and ambient witnesses}, with
$H=G_{\ge5}$ and $K=G_{>5}$.

\begin{theorem}[A colour with no local pair]
There is a colour $j\ge5$ such that every actual
$H$-orbit containing an ambient-eligible $j$-edge
and intransitive under $K$ has exactly one edge
eligible for its full local $j$-omission. Its
active labels form an interval starting at $5$,
and $\delta$ fixes it pointwise. Every $H$-orbit
contains at most one ambient-eligible $j$-edge,
and at least two distinct $H$-orbits contain one.
\end{theorem}
\begin{proof}
If no such colour existed, for each tail colour
one could choose a $K$-intransitive actual $H$-orbit
containing one ambient witness and two edges
eligible for the full local omission. These
independent choices would give exactly the mixed
assignment excluded by the complete acquired proof.
Thus a colour exists for which no such orbit
has a local pair.

Every $K$-intransitive witness orbit has at least
one locally eligible edge, namely its ambient
witness. Hence it has exactly one. Transitivity
of $H$ makes its local omission have exactly two
orbits, joined by that edge. If any $H$-orbit had
two ambient $j$-edges, then either it would already
be a forbidden $K$-intransitive target, or the
complete whole-$K$-orbit transport would carry
both edges simultaneously into such a target.
Both would remain ambient and therefore locally
eligible. This contradiction proves the global
at-most-one assertion.

The acquired active-label argument applies to
every witness orbit. A $K$-intransitive one starts
at $5$ and has a genuine local unique crossing.
The complete orbit-size and parity proof therefore
makes $\delta$ fix each of these orbits pointwise.
The two ambient edges required globally lie in
two distinct actual $H$-orbits. No transport of
the second edge, or extra local uniqueness for a
$K$-transitive orbit, is asserted here.
\end{proof}

\begin{lemma}[Prefixes of an anchored action]
Let $T=\gen{t_0,\ldots,t_{u-1}}$ be transitive on $V$
with string commutations. Suppose $a$ is fixed by
$t_1,\ldots,t_{u-1}$ and $t_{u-1}\ne1$.
For $1\le v<u$, put
$C_v=a^{\gen{t_0,\ldots,t_{v-1}}}$ and $C_0=\{a\}$.
Then every inclusion $C_{v-1}\subset C_v$ is strict,
so $|C_v|\ge v+1$. Every mark of index at least
$v+1$ fixes $C_v$ pointwise.
\end{lemma}
\begin{proof}
A later mark commutes with each generator of the
prefix and fixes $a$, proving the last assertion.
If $C_v=C_{v-1}$ with $v<u$, then $t_{v-1}$
preserves $C_{v-1}$; every mark of index at least
$v$ fixes $C_{v-1}$, by the same commutation
argument. Thus $C_{v-1}$ is invariant under all
of $T$. Transitivity makes it $V$, on which
$t_{u-1}$ would act trivially, a contradiction.
The strict inclusions and the degree bound follow.
We make no strict-growth claim at $v=u$.
\end{proof}

\begin{theorem}[A truncated local orbit]
Let $Q=\gen{\rho_m,\ldots,\rho_y}$ act transitively
on $W$, with string commutations and all displayed
restrictions nonidentity. Suppose its full local
$j$-omission has exactly two orbits $A,B$ joined
by exactly one $j$-edge $\{a,b\}$. For
$j+1\le x\le y$, at least $x-m+1$ points of $W$
are fixed by every mark of index greater than $x$.
\end{theorem}
\begin{proof}
Set $L=Q_{<j}$ and $R=Q_{>j}$. On each of $A,B$
the commuting product $LR$ is transitive. Every
mark at distance at least two from $j$ fixes
$a,b$, since it preserves the omission orbits
and the unique crossing edge.

Put $v=x-j-1\ge0$ and
$R'=\gen{\rho_{j+1},\ldots,\rho_{x-1}}$, with
$R'=1$ if $v=0$. The two sets
\[
                     C_A=a^{LR'},\qquad C_B=b^{LR'}
\]
are disjoint and fixed pointwise by all marks
after $x$: these marks fix both anchors and
commute with every displayed generator of $LR'$.
We bound their total size.

On each side the active labels of $L$ form an
interval ending at $j-1$, and those of $R$ an
interval starting at $j+1$, by the acquired
commuting anchored-factor argument. Write their
lengths as $h_A,u_A$ and $h_B,u_B$. Then
\[
 \max(h_A,h_B)=j-m=:h,\qquad \max(u_A,u_B)=y-j.
\]
All active left marks remain active on $C_A$ or
$C_B$ as appropriate: their own orbit through
the anchor is unchanged. The number of active
right marks on these sets is respectively
$\min(u_A,v)$ and $\min(u_B,v)$. If $v<u_A$,
this follows from the strict prefix lemma;
if $v\ge u_A$, it is the full right action.
The same argument applies to $B$. Their maximum
is $v$, since $v<y-j$.

Apply the complete two-commuting-factor size
bound to the transitive actions on $C_A,C_B$.
If the maxima $h,v$ occur on opposite sides,
the separate anchored bounds give total size
at least $h+v+2$. If they occur on one side,
that side gives $h+v+1$ and the other contributes
one more. Empty factors are included.

The only possible exception would have $h=v=1$,
$|C_A|=2$, $|C_B|=1$, after interchanging sides.
Then $L$ has just one active mark, acting as the
same transposition as $\rho_{j+1}$ on $C_A$.
Every further right mark fixes $a$ and commutes
with $L$, so fixes both points of $C_A$. It also
fixes $b$, while the first right mark and $L$
fix $b$ already. Thus $C_A,C_B$ are invariant
under the complete $LR$, forcing $A=C_A$ and
$B=C_B$. No right mark after $j+1$ could act
nontrivially on $W$, contradicting $y\ge x=j+2$.
The exception is impossible. Hence
\[
 |C_A|+|C_B|\ge h+v+2=x-m+1,
\]
as required. No assertion about two disjoint
point paths has entered the proof.
\end{proof}

\begin{theorem}[Two witnesses with intransitive first omissions]
For the colour of the first theorem, suppose two
distinct witness orbits $O_1,O_2$ are both
$K$-intransitive. Write their last active labels
as $x\le y$. Then $x<r-1$. If $x=y$,
\[
                   |\operatorname{Fix}(G_{>x})|\ge2x+1.
\]
If $j+1\le x<y$, the proved bound is
\[
                   |\operatorname{Fix}(G_{>x})|\ge2x.
\]
\end{theorem}
\begin{proof}
Both orbits are fixed pointwise by $\delta$ and
have first active label $5$. Their complete local
orbit bounds give $|O_1|\ge x-3$, $|O_2|\ge y-3$.
If $x=y$, both whole orbits and the seven head
points are fixed after $x$, giving $2x+1$.
If $j+1\le x<y$, the truncated theorem supplies
at least $x-4$ fixed points in $O_2$. Together
with $O_1$ and the seven head points, this gives
$(x-3)+(x-4)+7=2x$.

If $x=r-1$, then $y=r-1$ too. The two complete
orbit bounds and the seven head points occupy
at least $2(r-4)+7=2r-1$ of the $2r+1$ points.
The odd involution $\delta$ fixes this entire
union pointwise, and must move at least two
points outside it. Consequently exactly two
points remain and $\delta$ is literally their
transposition. The complete acquired exterior-pair
obstruction applies under the present actual
head and global ambient pair hypotheses. It
excludes this possibility, so $x<r-1$.
\end{proof}

\paragraph{Attribution and scope.}
The mixed-assignment contradiction, active intervals,
whole-orbit transport, complete orbit-size bound and
exterior-pair obstruction are acquired. The simultaneous
witness choice and strict anchored-prefix count supply
the new truncated bound. The unequal-end case currently
provides $2x$, not the $2x+1$ fixed points required for
the desired separator. The case $x=j<y$, a second
witness with transitive $K$-action, and global attachments
also remain. No complete L044 recovery, new induction,
computation, formal or question credit is claimed.
\end{document}
